\documentclass[11pt]{article}

\usepackage[margin=1in,headheight=22pt]{geometry}
\usepackage[T1]{fontenc}
\usepackage{amsmath}
\usepackage{newtxtext,newtxmath}
\usepackage{microtype}
\usepackage{abstract}
\usepackage{booktabs}
\usepackage{tabularx}
\usepackage{array}
\usepackage{caption}
\usepackage{enumitem}
\usepackage{fancyhdr}
\usepackage{titlesec}
\usepackage[table]{xcolor}
\usepackage{xurl}
\usepackage{hyperref}

\definecolor{ReviewBlue}{RGB}{31,78,121}
\definecolor{ReviewBlueLight}{RGB}{239,246,252}
\definecolor{ReviewRule}{RGB}{185,202,218}
\definecolor{ReviewGreen}{RGB}{221,237,221}
\definecolor{ReviewGreenLight}{RGB}{235,246,235}
\definecolor{ReviewYellow}{RGB}{255,243,205}
\definecolor{ReviewGray}{RGB}{238,238,238}

\hypersetup{
  colorlinks=true,
  linkcolor=ReviewBlue,
  citecolor=ReviewBlue,
  urlcolor=ReviewBlue,
  pdftitle={Art Gallery Problems for Orthogonal Polygons and Orthogonal Polyhedra},
  pdfsubject={Literature review},
  pdfauthor={},
  pdfkeywords={art gallery problem, orthogonal polygons, polyominoes, orthogonal polyhedra, perimeter bounds}
}

\setcounter{tocdepth}{2}
\setlength{\parindent}{1.2em}
\setlength{\parskip}{0pt}

\renewcommand{\abstractnamefont}{\normalfont\large\bfseries\color{ReviewBlue}}
\renewcommand{\abstracttextfont}{\normalfont\small}
\setlength{\absleftindent}{0.07\textwidth}
\setlength{\absrightindent}{0.07\textwidth}

\captionsetup{
  font=small,
  labelfont={bf,color=ReviewBlue},
  textfont=it,
  width=0.94\textwidth
}

\titleformat{\section}
  {\normalfont\Large\bfseries\color{ReviewBlue}}
  {\thesection}{0.75em}{}
\titleformat{\subsection}
  {\normalfont\large\bfseries}
  {\thesubsection}{0.65em}{}
\titleformat{\subsubsection}
  {\normalfont\normalsize\bfseries}
  {\thesubsubsection}{0.55em}{}
\titlespacing*{\section}{0pt}{2.1ex plus 0.6ex minus 0.2ex}{1.0ex plus 0.2ex}
\titlespacing*{\subsection}{0pt}{1.6ex plus 0.4ex minus 0.2ex}{0.7ex plus 0.2ex}

\pagestyle{fancy}
\fancyhf{}
\lhead{\small\textsc{Orthogonal Art Gallery Problems}}
\rhead{\small\thepage}
\renewcommand{\headrulewidth}{0.4pt}
\renewcommand{\headrule}{\hbox to\headwidth{\color{ReviewRule}\leaders\hrule height \headrulewidth\hfill}}
\fancypagestyle{plain}{%
  \fancyhf{}
  \rhead{\small\thepage}
  \renewcommand{\headrulewidth}{0pt}
}

\setlist[itemize]{leftmargin=*, itemsep=0.35em}
\setlist[enumerate]{leftmargin=*, itemsep=0.45em}
\setlist[description]{style=nextline, leftmargin=1.5em, labelsep=0.5em, itemsep=0.45em, font=\normalfont\bfseries}

\newcolumntype{Y}{>{\raggedright\arraybackslash}X}
\newcommand{\floor}[1]{\left\lfloor #1 \right\rfloor}
\newcommand{\ceil}[1]{\left\lceil #1 \right\rceil}
\newcommand{\loversix}{\floor{\ell/6}}
\newcommand{\covered}[1]{\cellcolor{ReviewGreen}#1}
\newcommand{\weaker}[1]{\cellcolor{ReviewGreenLight}#1}
\newcommand{\limited}[1]{\cellcolor{ReviewBlueLight}#1}
\newcommand{\uncertaincell}[1]{\cellcolor{ReviewYellow}#1}
\newcommand{\outside}{\cellcolor{ReviewGray}--}
\newcommand{\reviewcallout}[2]{%
  \begin{center}
  \setlength{\fboxsep}{8pt}
  \fcolorbox{ReviewBlue}{ReviewBlueLight}{%
  \begin{minipage}{0.91\textwidth}
  {\bfseries\color{ReviewBlue}#1}\par\vspace{0.35em}
  #2
  \end{minipage}}
  \end{center}
}

\begin{document}
\thispagestyle{plain}
\begin{center}
{\LARGE\bfseries Art Gallery Problems for Orthogonal Polygons and Orthogonal Polyhedra\par}
\vspace{0.5em}
{\large A Literature Review of Results and Open Problems\par}
\vspace{0.85em}
{\normalsize May 23, 2026\par}
\vspace{1.1em}
\textcolor{ReviewRule}{\rule{0.88\textwidth}{0.6pt}}
\end{center}

\begin{abstract}
This review surveys art-gallery problems for orthogonal polygons,
polyominoes, and orthogonal polyhedra, with emphasis on perimeter-type bounds,
holes, and the model changes that produce current open problems.  In two
dimensions, the perimeter-over-six question is meaningful only after fixing a
lattice scale.  Ma{\ss}berg proves the tight $\floor{\ell/6}$ point-guard
bound for hole-free polyominoes, and that theorem also covers hole-free
integral orthogonal polygons after unit-grid subdivision.  The unresolved
perimeter-over-six direction is therefore the version with holes, or a
broader lattice-domain formulation not reducible to the hole-free polyomino
case.  In discrete 3D and higher dimensions, the cell-volume analogue is known:
Pinciu proves the tight $\floor{(m+1)/3}$ point-guard theorem for
$m$-polyhypercubes.  For arbitrary continuous orthogonal polyhedra, however,
point guards do not give a linear planar analogue: some examples require
$\Theta(n^{3/2})$ guards.  The continuous 3D literature consequently turns to
edge guards, reflex-edge guards, and face guards, while discrete voxel/cell
guard variants form a separate line of open problems.
\end{abstract}

\tableofcontents
\clearpage

\section{Executive Synthesis}

The review has one organizing question: which art-gallery bounds survive when
the object is orthogonal, the perimeter is lattice-normalized, holes may be
present, and the dimension changes?  The answer is not a single theorem but a
map of comparable models.  The sharp perimeter-over-six theorem is known for
hole-free polyominoes, and the same theorem settles the hole-free integral
orthogonal-polygon reading by subdivision.  The remaining two-dimensional
frontier is the holes version, together with any broader lattice-domain
formulation that does not reduce to a hole-free polyomino.  In discrete 3D and
higher dimensions, Pinciu gives the direct point-guard cell-volume theorem for
polyhypercubes.  In continuous orthogonal polyhedra, the direct point-guard
analogy fails, and edge, reflex-edge, face, and cell-guard variants carry the
relevant open questions.

\begin{table}[htbp]
\small
\setlength{\tabcolsep}{4pt}
\renewcommand{\arraystretch}{1.18}
\begin{tabularx}{\textwidth}{
  >{\raggedright\arraybackslash}p{0.25\textwidth}
  >{\raggedright\arraybackslash}p{0.22\textwidth}
  Y
  >{\raggedright\arraybackslash}p{0.18\textwidth}}
\toprule
Result line & Current status & Why it matters here & Main evidence \\
\midrule
Simple orthogonal polygons & Tight $\floor{n/4}$ vertex/point-guard baseline
for hole-free orthogonal polygons. & Establishes the classical 2D orthogonal
benchmark, but it is a vertex-count theorem, not a perimeter theorem. &
Kahn--Klawe--Kleitman; Gyori
\cite{KahnKlaweKleitman1983,Gyori1986} \\

Hole-free polyomino perimeter & Tight
$\max\{1,\floor{\ell/6}\}$ point-guard theorem. & This is the core
perimeter-over-six result and the source of the hole-free integral-polygon
overlap. & Ma{\ss}berg \cite[Theorem~1]{Massberg2014} \\

Integral orthogonal polygons, no holes & Solved by reduction to hole-free
polyominoes after unit-grid subdivision. & Prevents the Diaz-Banez et al.
$\floor{N/6}$ question from being misread as open in the hole-free case. &
Ma{\ss}berg; Diaz-Banez et al.
\cite[Theorem~1]{Massberg2014}; \cite{DiazBanezEtAl2025} \\

Polyomino or integral domains with holes & Open or conditional for
perimeter-over-six point guards. & This is the main 2D open direction closest
to the original perimeter question. & Ma{\ss}berg habilitation
\cite[Conj.~6.10]{MassbergHabilitation} \\

Orthogonal polygons with holes, vertex guards & Partially solved; strongest
conjectured bounds remain open. & These results explain the role of holes, but
they do not settle point-guard perimeter bounds. & Urrutia; Zylinski;
Michael--Pinciu
\cite{UrrutiaOpenProblems,Zylinski2006,MichaelPinciu2016} \\

Polycubes and polyhypercubes & Tight $\floor{(m+1)/3}$ point-guard
cell-volume theorem. & This is the genuine 3D-and-higher analogue of the
polyomino cell-count theorem, but only for unit-cell complexes. & Pinciu
\cite{Pinciu2015} \\

Continuous orthogonal polyhedra in 3D & No direct point-guard analogue; edge
and reflex-edge variants remain active. & The continuous 3D branch must be
read through guard model changes rather than by copying the planar perimeter
theorem. &
Paterson--Yao; Viglietta; Benbernou et al.; Cano et al.
\cite{PatersonYao1992,VigliettaThesis,BenbernouEtAl2011,CanoTothUrrutiaViglietta2022} \\
\bottomrule
\end{tabularx}
\caption{Reader-facing synthesis of the central result lines.  Each row is
intended to be read only under the stated model; later sections give the
definitions, non-implications, and evidence trail.}
\label{tab:executive-synthesis}
\end{table}

The document is assembled in layers.  The front matter gives the short answer;
the terminology section fixes the model axes; the historical flow explains how
the result lines diverged; the status matrix and coverage chart show what is
proved, conditional, or open; the non-implication and proof-technique sections
explain why adjacent theorems do not automatically transfer; the model-by-model
sections give the detailed survey; and the open-problem, source-guide, and
evidence sections make the unresolved directions auditable.

\section{Introduction and Scope}

The art-gallery literature is best read through four modeling choices: the
object class, the presence of holes, the allowed guard model, and the
parameter in the bound.  A statement about point guards in a hole-free
polyomino need not imply a statement about vertex guards in a polygon with
holes, and neither statement transfers automatically to edge guards in an
orthogonal polyhedron.  This review keeps those distinctions explicit because
most apparent conflicts among the results come from changing one of these
choices.

For two-dimensional orthogonal objects, the main perimeter issue is scale.
There is no scale-free theorem of the form ``at most $P/6$ guards'' for
arbitrary real-coordinate orthogonal polygons when $P$ is Euclidean
perimeter: scaling the polygon changes $P$ without changing the guard number.
Perimeter-over-six results therefore belong to lattice-normalized settings,
such as polyomino perimeter, integral perimeter, or the unit-edge
ortho-unit model.

\reviewcallout{Main 2D Distinction}{%
If $Q$ is a hole-free integral orthogonal polygon of lattice perimeter $N$,
unit-grid subdivision identifies it with a hole-free polyomino of perimeter
$\ell=N$.  Ma{\ss}berg's theorem then gives
$\max\{1,\floor{N/6}\}$ point guards, and hence $\floor{N/6}$ guards whenever
$N\geq 6$.  Thus the nontrivial hole-free integral-polygon reading of the
Diaz-Banez et al. $\floor{N/6}$ question is already covered by Ma{\ss}berg;
the genuinely unresolved reading is the extension to holes or to a broader
lattice-domain convention.}

In three dimensions the direct point-guard analogy breaks for a different
reason.  Orthogonal polyhedra can require superlinear numbers of point guards
in the number of vertices, so the 3D branch of the literature usually asks for
edge guards, reflex-edge guards, face guards, or guards in discrete polycube
visibility models.  The word \emph{orthonormal} is not the relevant geometric
term; the standard objects are orthogonal polygons, orthogonal polyhedra, and
orthogonal polytopes.

The following table summarizes the readings of the perimeter-over-six phrase
that recur in the literature.

\begin{center}
\small
\setlength{\tabcolsep}{5pt}
\renewcommand{\arraystretch}{1.15}
\begin{tabularx}{0.96\textwidth}{>{\raggedright\arraybackslash}p{0.34\textwidth} Y Y}
\toprule
Reading of the $\floor{N/6}$ question & Status & Reason \\
\midrule
Simple/hole-free integral orthogonal polygon & Solved & Same object as a
hole-free polyomino after unit-grid subdivision; $N=\ell$. \\
Ortho-unit polygon & Solved, stronger & Diaz-Banez et al. prove
$\floor{(n-4)/8}$, and here $n=P=N$. \\
Integral or polyomino domain with holes & Open/conditional & Ma{\ss}berg's
hole extension depends on an unproved maximal-rectangle packing conjecture. \\
Real-coordinate orthogonal polygon with Euclidean perimeter & Not a valid
scale-free statement & Uniform scaling changes perimeter but not guard number. \\
\bottomrule
\end{tabularx}
\end{center}

\section{Conventions and Terminology}

\paragraph{2D objects.}
An \emph{orthogonal polygon} is a polygon whose edges are horizontal or
vertical; the literature also uses \emph{rectilinear} for the same class.  A
\emph{polyomino} is a finite edge-connected union of unit axis-parallel
squares; its perimeter is the number of unit boundary edges, including hole
boundaries when holes are present.  An \emph{integral orthogonal polygon} has
integer coordinates, so its perimeter is measured as a lattice length.  An
\emph{ortho-unit polygon} is an orthogonal polygon all of whose edges have
length one; in that model, the number of vertices equals the perimeter.

\paragraph{3D objects.}
The standard 3D term is \emph{orthogonal polyhedron}: a polyhedron whose faces
are axis-aligned, equivalently whose face normals are coordinate directions.
In higher dimensions, the corresponding term is \emph{orthogonal polytope}.
The word \emph{orthonormal} describes vectors or bases, not galleries.
Polycubes are the voxel analogue of polyominoes and form a discrete subclass
of orthogonal polyhedra.  A polyhypercube is the corresponding finite
face-connected union of unit hypercubes in dimension $d$.

\paragraph{Guard models.}
Unless a section says otherwise, a 2D guard is a standard point guard with
straight-line visibility inside the polygon.  In 3D the model must be stated:
the literature distinguishes point guards, vertex guards, edge guards,
reflex-edge guards, face guards, open edge guards, and discrete polycube
visibility graph models.  These models are not interchangeable.

\paragraph{Why perimeter needs a lattice scale.}
No theorem of the form ``at most $P/6$ guards'' can hold for arbitrary
real-coordinate orthogonal polygons if $P$ is Euclidean perimeter: scaling a
polygon down preserves the guard number while making $P$ arbitrarily small.
The perimeter-over-six question is therefore a lattice-normalized question:
polyomino perimeter, integral perimeter, or a fixed unit-edge model.

\paragraph{Common parameters.}
The same letters are used differently across the literature, so the review
uses the following conventions unless a cited theorem uses its own notation.

\begin{center}
\small
\setlength{\tabcolsep}{5pt}
\renewcommand{\arraystretch}{1.12}
\begin{tabularx}{0.94\textwidth}{>{\raggedright\arraybackslash}p{0.12\textwidth} Y Y}
\toprule
Symbol & Meaning in this review & Main caution \\
\midrule
$n$ & Number of polygon vertices, or polyhedron vertices when the 3D context is
explicit. & Vertex-count bounds are not perimeter bounds. \\
$h$ & Number of holes in a 2D polygonal domain. & Holes can change both
decomposition structure and guard type. \\
$r$ & Number of reflex edges in a 3D orthogonal polyhedron. & This is a 3D
parameter; in 2D, reflex vertices play the analogous role. \\
$m$ & Number of polyomino cells in 2D, or number of polyhedron edges in 3D
when stated. & Area/cell bounds and edge-count bounds should not be compared
without the model. \\
$\ell$ & Lattice perimeter of a polyomino, including hole boundaries. & This is
the parameter in Ma{\ss}berg's perimeter theorem. \\
$N$ & Lattice perimeter of an integral orthogonal polygon. & After unit-grid
subdivision, $N=\ell$ for the corresponding polyomino. \\
$P$ & Euclidean perimeter when discussing scale, or a polygon only when no
formula is involved. & Real-coordinate Euclidean perimeter is not scale-free. \\
\bottomrule
\end{tabularx}
\end{center}

\section{Historical Flow}

\begin{enumerate}
  \item Classical art-gallery theorems start with arbitrary simple polygons:
  Chvatal proves $\floor{n/3}$ and Fisk gives the standard triangulation proof
  \cite{Chvatal1975,Fisk1978}.  Shermer, O'Rourke, and Urrutia provide the
  main survey/book context \cite{Shermer1992,ORourke1987,Urrutia2000}.

  \item Orthogonal polygons admit stronger vertex-count theorems:
  Kahn--Klawe--Kleitman prove the tight $\floor{n/4}$ bound for simple
  orthogonal polygons \cite{KahnKlaweKleitman1983}.

  \item Holes complicate the 2D theory.  Vertex-guard conjectures for
  orthogonal polygons with holes, especially Shermer's and Hoffmann's
  formulations, remain separate from the perimeter-over-six point-guard
  question \cite{UrrutiaOpenProblems,HoffmannKriegel1996,Zylinski2006,MichaelPinciu2016}.

  \item Polyominoes introduce area and lattice-perimeter parameters.  The
  tight cell-count theorem covers connected $m$-polyominoes with $m\geq 2$,
  holes allowed \cite[Cor.~1]{BiedlEtAl2012}, while Ma{\ss}berg's
  perimeter-over-six theorem is proven for hole-free polyominoes and becomes
  conditional in the presence of holes
  \cite[Theorem~1]{Massberg2014}; \cite[Thm.~6.3 and
  Conj.~6.10]{MassbergHabilitation}.

  \item Recent ortho-unit work proves a strong unit-edge theorem and states
  an integral-perimeter conjecture \cite{DiazBanezEtAl2025}.  Its literal
  hole-free integral-polygon reading overlaps with Ma{\ss}berg's polyomino
  theorem; only the holes or broader-domain reading remains open.

  \item In 3D, point guards are too weak for a planar-style theorem:
  orthogonal polyhedra need $\Theta(n^{3/2})$ point guards in the worst case
  \cite{PatersonYao1992,VigliettaThesis}.  Edge-guard and reflex-edge-guard
  conjectures become the main 3D analogue \cite{BenbernouEtAl2011,Viglietta2020,CanoTothUrrutiaViglietta2022}.
\end{enumerate}

\section{Status Overview}

Table~\ref{tab:status-matrix} separates the main models discussed below.  The
labels ``proven tight,'' ``conditional,'' ``conjectural,'' and ``open'' refer
to the exact model in the first three columns; changing the guard type,
allowing holes, or changing the parameter can change the answer.

\begin{table}[htbp]
\small
\setlength{\tabcolsep}{3.5pt}
\renewcommand{\arraystretch}{1.16}
\begin{tabularx}{\textwidth}{
  >{\raggedright\arraybackslash}p{0.18\textwidth}
  >{\raggedright\arraybackslash}p{0.11\textwidth}
  >{\raggedright\arraybackslash}p{0.13\textwidth}
  Y
  >{\raggedright\arraybackslash}p{0.17\textwidth}}
\toprule
Model & Guards & Parameter & Status & Main source \\
\midrule
Simple orthogonal polygon & point & vertices $n$ &
Proven tight: $\floor{n/4}$ & Kahn--Klawe--Kleitman; Gyori
\cite{KahnKlaweKleitman1983,Gyori1986} \\

Orthogonal polygon with holes & vertex & vertices $n$, holes $h$ &
Open at the strongest conjectured bounds; improved h-independent upper bounds
and partial sharp cases are known & Urrutia; Hoffmann--Kriegel; Zylinski;
Michael--Pinciu
\cite{UrrutiaOpenProblems,HoffmannKriegel1996,Zylinski2006,MichaelPinciu2016} \\

Hole-free polyomino & point & lattice perimeter $\ell$ &
Proven tight for $\ell\geq 6$: $\loversix$ & Ma{\ss}berg
\cite[Theorem~1]{Massberg2014} \\

Polyomino with holes & point & lattice perimeter $\ell$ &
Conditional/open: Ma{\ss}berg's extension depends on a rectangle-packing
conjecture & Ma{\ss}berg habilitation
\cite[Conj.~6.10]{MassbergHabilitation} \\

Polyomino with holes & point & cells $m$ &
Proven tight for $m\geq 2$: $\floor{(m+1)/3}$; $m=1$ is trivial; does not
imply perimeter-over-six & Biedl et al. \cite[Cor.~1]{BiedlEtAl2012} \\

Polyhypercube & point & cells $m$ &
Proven tight for $m\geq 2$ and $d\geq 2$: $\floor{(m+1)/3}$; this is the
cell-volume analogue, not a continuous polyhedron theorem & Pinciu
\cite{Pinciu2015} \\

Polyhypercube & voxel/cell & cells $m$, dimension $d$ &
Open with partial bounds: Pinciu gives lower and upper bounds and conjectures
a sharp dependent-on-$d$ formula & Pinciu \cite{Pinciu2015} \\

Ortho-unit polygon & point & perimeter $P=n$ &
Proven tight for $n\geq 12$: $\floor{(n-4)/8}$, stronger than $P/6$ in this
model & Diaz-Banez et al. \cite{DiazBanezEtAl2025} \\

Integral orthogonal polygon, hole-free & point & lattice perimeter $N$ &
Solved by overlap: subdivision gives a hole-free polyomino with $N=\ell$, so
Ma{\ss}berg proves the nontrivial Diaz-Banez $\floor{N/6}$ bound in this case &
Ma{\ss}berg; Diaz-Banez et al.
\cite[Theorem~1]{Massberg2014}; \cite{DiazBanezEtAl2025} \\

Integral or polyomino domain with holes & point & lattice perimeter $N$ or
$\ell$ & Still open/conditional: Ma{\ss}berg's hole extension depends on a
rectangle-packing conjecture & Ma{\ss}berg habilitation
\cite[Conj.~6.10]{MassbergHabilitation} \\

Orthogonal polyhedron & point & vertices $n$ &
No linear planar analogue: worst case is $\Theta(n^{3/2})$ & Paterson--Yao;
Viglietta \cite{PatersonYao1992,VigliettaThesis} \\

Orthogonal polyhedron & edge or reflex edge & edges $m$, reflex edges $r$ &
Main 3D analogue; Urrutia-type conjectures remain open beyond partial classes
& Benbernou et al.; Viglietta; Cano et al.
\cite{BenbernouEtAl2011,Viglietta2020,CanoTothUrrutiaViglietta2022} \\
\bottomrule
\end{tabularx}
\caption{Status map for orthogonal art-gallery problems closest to
perimeter-over-six.}
\label{tab:status-matrix}
\end{table}

Figure~\ref{fig:2d-coverage} summarizes which two-dimensional polygon classes
are actually covered by the main bounds in this review.
Green cells are proven coverage, pale green cells are proven but usually
weaker than a later special theorem, blue cells indicate a proper-subclass or
different-parameter theorem, yellow cells mark conditional or open status, and
gray cells are outside the theorem as stated.

\begin{figure}[htbp]
\centering
\scriptsize
\setlength{\tabcolsep}{2pt}
\renewcommand{\arraystretch}{1.22}
\begin{tabularx}{\textwidth}{
  >{\raggedright\arraybackslash}p{0.18\textwidth}
  *{6}{>{\centering\arraybackslash}X}}
\toprule
2D object class &
Classical point $\floor{n/3}$ &
Orthogonal point $\floor{n/4}$ &
Orthogonal holes, vertex guards &
Perimeter point $\floor{\ell/6}$ or $\floor{N/6}$ &
Area point $\floor{(m+1)/3}$ &
Ortho-unit point $\floor{(n-4)/8}$ \\
\midrule
General simple polygon, no holes &
\covered{yes, tight} &
\outside &
\outside &
\outside &
\outside &
\outside \\

Simple orthogonal polygon, no holes &
\weaker{yes, weaker} &
\covered{yes, tight} &
\outside &
\outside &
\outside &
\outside \\

Hole-free integral orthogonal polygon &
\weaker{yes, weaker} &
\weaker{yes, vertex bound} &
\outside &
\covered{yes, by subdivision} &
\limited{yes, area parameter} &
\limited{only unit-edge subclass} \\

Hole-free polyomino &
\weaker{yes, weaker} &
\weaker{yes, vertex bound} &
\outside &
\covered{yes, tight} &
\limited{yes, area parameter} &
\limited{only unit-edge subclass} \\

Ortho-unit polygon &
\weaker{yes, weaker} &
\weaker{yes, vertex bound} &
\outside &
\covered{yes, follows} &
\limited{yes, area parameter} &
\covered{yes, tight} \\

Orthogonal domain with holes, real coordinates &
\outside &
\outside &
\uncertaincell{partial/open} &
\outside &
\outside &
\outside \\

Integral or polyomino domain with holes &
\outside &
\outside &
\uncertaincell{vertex-only partial/open} &
\uncertaincell{cond./open} &
\covered{yes, tight} &
\outside \\
\bottomrule
\end{tabularx}
\caption{Two-dimensional coverage chart for the principal bounds.  A dash
means the theorem does not apply as stated, usually because the object has
holes, lacks a lattice scale, uses the wrong guard model, or is controlled by a
different parameter.}
\label{fig:2d-coverage}
\end{figure}

\clearpage

\section{Model Separations and Non-Implications}

The most useful way to test whether a proposed theorem is genuinely new is to
ask which model changes are being made.  The following separations account for
most apparent overlaps in this part of the literature.

\begin{description}
  \item[Scale separation.]
  A Euclidean perimeter statement for arbitrary real-coordinate orthogonal
  polygons cannot be scale-free.  Any perimeter-over-six theorem must specify a
  unit: polyomino boundary length, integral lattice length, or unit-edge
  ortho-unit length.

  \item[Point guards versus vertex guards.]
  Vertex-guard theorems for orthogonal polygons with holes do not imply
  point-guard perimeter theorems.  Conversely, a point-guard perimeter theorem
  need not place guards on vertices.  This is why Shermer's and Hoffmann's
  vertex-guard conjectures are adjacent to, but distinct from, the polyomino
  perimeter question.

  \item[Standard visibility versus restricted visibility.]
  Floodlights, half-plane guards, sliding cameras, transmitters, $k$-hop
  visibility, and rook/queen polycube models are adjacent to standard
  point-guard or vertex-guard visibility, but they are not interchangeable
  with it.  They can supply techniques and open problems, but each requires a
  separate model cell before its theorems can be used in this review.  The
  current evidence matrix has separate cells for the currently cited
  sliding-camera, sliding-transmitter, $k$-hop, polyhypercube point-guard,
  polyhypercube pixel/voxel-guard, face-guard, and 2-reflex variants.  It
  also has explicit scope-exclusion cells for
  floodlight, half-plane, dispersive, contiguous, mobile, point-boundary, and
  related variants that have not been promoted into theorem-level discussion.

  \item[Holes versus hole-free domains.]
  The hole-free perimeter theorem is not automatically stable under adding
  holes.  Holes change the intersection graph of maximal rectangles and break
  the perfect-graph route used in the hole-free proof.  Ma{\ss}berg's
  maximal-rectangle packing conjecture is precisely the missing replacement
  structure in that setting.

  \item[Area versus perimeter.]
  The tight $\floor{(m+1)/3}$ theorem for connected $m$-cell polyominoes with
  holes, $m\geq 2$, is an area theorem; the one-cell case is the trivial
  one-guard case.  It does not imply a perimeter theorem because polyominoes
  can have very different area-to-perimeter ratios, and the conjectured
  perimeter bound is sensitive to boundary complexity rather than cell count.

  \item[Ortho-unit versus integral.]
  Ortho-unit polygons form a narrow unit-edge subclass: every edge has length
  one, so vertex count and perimeter coincide.  Integral orthogonal polygons
  allow longer lattice edges.  In the hole-free case, subdivision reduces the
  integral model to hole-free polyominoes; with holes, it reduces to the
  polyomino-with-holes case, which is still conditional/open for perimeter.

  \item[2D point guards versus 3D point guards.]
  Orthogonal polyhedra do not inherit planar point-guard behavior.  The
  $\Theta(n^{3/2})$ point-guard bound in 3D forces the analogue to move toward
  edge guards, reflex-edge guards, face guards, or discrete polycube models.
\end{description}

\section{Proof Techniques and Where They Break}

\begin{description}
  \item[Triangulation and coloring.]
  Chvatal and Fisk establish the classical template for simple polygons
  \cite{Chvatal1975,Fisk1978}.  Kahn--Klawe--Kleitman and Gyori adapt
  coloring ideas to simple orthogonal polygons
  \cite{KahnKlaweKleitman1983,Gyori1986}.  The method is powerful for
  vertex-count bounds but does not directly produce a lattice-perimeter
  theorem.

  \item[Quadrilateralization and cactus duals.]
  Orthogonal holes are often attacked by decomposing the domain into
  rectangles or quadrilaterals and coloring a dual structure.  Zylinski's
  cactus-dual condition explains why the $\floor{(n+h)/4}$ conjecture is
  approachable in some cases but not yet resolved in full
  \cite{Zylinski2006}.  Michael and Pinciu introduce same-sign diagonal graphs
  of convex quadrangulations and use vertex-cover bounds to unify several
  orthogonal-guarding results and improve h-independent bounds for polygons
  with holes \cite{MichaelPinciu2016}.

  \item[Rectangle packing and perfect-graph methods.]
  Motwani--Raghunathan--Saran, Worman--Keil, and Ma{\ss}berg connect
  orthogonal guarding with star polygons, rectangle systems, and perfect
  graphs \cite{MotwaniRaghunathanSaran1990,WormanKeil2007,Massberg2014}.
  This is the proof technology behind the hole-free polyomino perimeter
  theorem.  Its failure point is explicit: the hole case depends on
  Ma{\ss}berg's maximal-rectangle packing conjecture
  \cite[Conj.~6.10]{MassbergHabilitation}.

  \item[Area and grid-cell arguments.]
  Biedl et al. prove the tight $\floor{(m+1)/3}$ theorem for connected
  $m$-polyominoes with $m\geq 2$ by using the cell structure rather than
  perimeter; the sufficiency proof works even for restrictive $r$-visibility
  \cite[Theorem~1 and Cor.~1]{BiedlEtAl2012}.  This is strong evidence that
  holes are tractable in a discrete model, but it leaves open whether
  perimeter alone controls the guard number.

  \item[Hardness and approximation reductions.]
  Schuchardt--Hecker, Katz--Roisman, and later work on sliding cameras and
  transmitters show that optimization versions remain hard even in restricted
  orthogonal settings
  \cite{SchuchardtHecker1995,KatzRoisman2008,DurocherEtAl2017,BiedlEtAl2019}.
  These results do not contradict extremal upper bounds; they warn that
  finding a minimum guard set is a different problem from proving a universal
  bound.

  \item[3D decomposition and binary-space partitions.]
  Paterson--Yao's binary-space-partition theorem, as used in Viglietta's
  survey, shows that point guards in orthogonal polyhedra have
  $\Theta(n^{3/2})$ worst-case behavior
  \cite{PatersonYao1992,VigliettaThesis}.  This is the main reason a 3D
  analogue should not simply ask for a point-guard version of the planar
  perimeter theorem.

  \item[Edge-guard and reflex-edge induction.]
  Benbernou et al., Viglietta, and Cano--Toth--Urrutia--Viglietta show that
  the most promising 3D analogue is based on edges, reflex edges, and genus
  rather than point guards
  \cite{BenbernouEtAl2011,Viglietta2020,CanoTothUrrutiaViglietta2022}.  The
  remaining gap is structural: the known upper bounds do not meet the
  Urrutia-type lower-bound constructions.
\end{description}

\section{Proof Roadmap for Core Result Lines}

This section records what each main proof line actually buys.  It is included
to make the survey useful for readers who want to enter an open problem rather
than only know its status.

\subsection{Coloring Proofs in the Plane}

The Chvatal--Fisk line proves existence by decomposing a polygon into simple
pieces, coloring a graph derived from that decomposition, and choosing the
smallest color class \cite{Chvatal1975,Fisk1978}.  The orthogonal
Kahn--Klawe--Kleitman theorem uses the extra parity and degree structure of
rectilinear polygons to improve the general $\floor{n/3}$ bound to
$\floor{n/4}$ \cite{KahnKlaweKleitman1983}; Gyori's proof is the shortest
entry point for this result \cite{Gyori1986}.  The transferable idea is the
small-color-class argument.  The non-transferable part is the parameter: these
proofs count vertices, not lattice perimeter.

\subsection{Holes and Quadrilateralization}

For orthogonal polygons with holes, the natural decomposition is a
quadrilateralization rather than a triangulation.  The dual structure can fail
to be a tree, so the coloring argument needs additional hypotheses.  Zylinski
isolates one such hypothesis through cactus-dual quadrilateralizations and
recovers Aggarwal's theorem for $h \leq 2$ holes
\cite{Zylinski2006}.  Michael--Pinciu give a different organizing language:
same-sign diagonal graphs and vertex covers \cite{MichaelPinciu2016}.  These
papers show that holes are not merely a nuisance parameter; they alter the
combinatorial object being colored or covered.

This is why the hole literature should be cited carefully.  The results are
highly relevant to orthogonal domains with holes, but their strongest
statements are usually vertex-guard and vertex-count statements.  They do not
settle point-guard perimeter-over-six.

\subsection{Perfect Graphs, Rectangles, and Perimeter}

Ma{\ss}berg's perimeter theorem sits in the rectangle/perfect-graph branch of
the literature rather than the triangulation/coloring branch.  The proof
relates guarding to systems of maximal rectangles, uses perfect-graph
structure to control the relevant packing/covering quantity, and then charges
that structure to the lattice perimeter \cite[Section~4]{Massberg2014}.  This
is why the theorem has the form $\max\{1,\floor{\ell/6}\}$ and why it is
genuinely different from the $\floor{n/4}$ orthogonal-polygon theorem.

The same roadmap also explains the open hole case.  In a rectilinear gallery
with holes, the intersection graph of maximal rectangles need not retain the
same perfect-graph behavior.  Ma{\ss}berg's habilitation identifies a
maximal-rectangle packing conjecture as the missing structural substitute:
for rectilinear galleries with holes, the maximum packing size of maximal
rectangles should upper-bound the number of guards required
\cite[Conj.~6.10]{MassbergHabilitation}.  A proof of that conjecture, combined
with the perimeter charging lemma for polyominoes with possible holes, would
prove the perimeter theorem with holes \cite[Lem.~6.9]{MassbergHabilitation};
a counterexample would likely force either a correction term or a different
invariant.

\subsection{Cell-Count Polyomino Bounds}

Biedl et al. prove the tight $\floor{(m+1)/3}$ theorem for connected
$m$-cell polyominoes with $m\geq 2$, including holes
\cite[Cor.~1]{BiedlEtAl2012}.  The lower bound is given by a comb family, and
the upper bound follows from a constructive decomposition that also works
under restrictive $r$-visibility \cite[Fig.~3 and Theorem~1]{BiedlEtAl2012}.
The one-cell case is handled separately by the obvious one guard.  This proof
line uses the cell structure directly, which is why it survives holes cleanly.
The result is strong, but it answers a different extremal question: how guard
number scales with area, not with boundary complexity.  For research purposes,
the important lesson is that holes are not inherently fatal in polyomino
models; they are fatal to the specific perimeter proof technology currently
known.

\subsection{3D Decomposition and Guard Models}

In three dimensions, the first modeling decision is guard type.  The
Paterson--Yao/Viglietta point-guard line shows that orthogonal polyhedra can
require $\Theta(n^{3/2})$ point guards \cite{PatersonYao1992,VigliettaThesis}.
That fact rules out a direct planar-style point-guard theorem in terms of
vertices or edges.

The promising 3D analogy therefore uses edges.  Benbernou et al. prove
open-edge-guard bounds for orthogonal polyhedra, including bounds expressed in
terms of total edges and reflex-edge count
\cite[Theorem~5 and Corollaries~6--7]{BenbernouEtAl2011}; Viglietta proves
tight-style results for 2-reflex orthogonal polyhedra \cite{Viglietta2020}; and
Cano--Toth--Urrutia--Viglietta improve the general polyhedron edge-guard
bound \cite{CanoTothUrrutiaViglietta2022}.  The research gap is not whether
edge guards are relevant, but whether the existing decomposition and charging
methods can be sharpened to the lower-bound scale.

\section{Status by Model}

\begin{description}
  \item[Simple orthogonal polygons, point guards.]
  Proven tight: every $n$-vertex simple orthogonal polygon is guardable by
  $\floor{n/4}$ point guards \cite{KahnKlaweKleitman1983,Gyori1986}.

  \item[Orthogonal polygons with holes, vertex guards.]
  Several upper bounds and partial sharp results are known.  O'Rourke gives
  the classical $\floor{(n+2h)/4}$ style bound for $h$ holes
  \cite{ORourke1987}.  Shermer's stronger $\floor{(n+h)/4}$ conjecture is
  open in general; Zylinski gives a coloring proof of Aggarwal's theorem for
  $h \leq 2$ and cactus-dual quadrilateralizations
  \cite{UrrutiaOpenProblems,Zylinski2006}.  In the h-independent direction,
  Hoffmann and Kriegel prove an $n/3$ upper bound, and Michael--Pinciu improve
  this line to $\floor{(17n-8)/52}$ using diagonal graphs and vertex covers
  \cite{HoffmannKriegel1996,MichaelPinciu2016}.

  \item[Hole-free polyominoes, point guards, lattice perimeter $\ell$.]
  Proven tight: for $\ell \geq 6$, Ma{\ss}berg proves that $\loversix$ point
  guards suffice \cite[Theorem~1]{Massberg2014}.  The exact theorem is
  $\max\{1,\loversix\}$, which includes the one-square case.

  \item[Polyominoes with holes, point guards, lattice perimeter $\ell$.]
  Open/conditional.  Ma{\ss}berg's text explicitly states that the extension
  to holes would follow from a maximal-rectangle packing conjecture
  \cite[Lem.~6.9 and Conj.~6.10]{MassbergHabilitation}.

  \item[Polyominoes with holes, point guards, area $m$.]
  Proven tight: every connected $m$-cell polyomino with $m\geq 2$, holes
  allowed, is guardable by $\floor{(m+1)/3}$ point guards
  \cite[Cor.~1]{BiedlEtAl2012}.  The one-cell case needs one guard.  This
  does not settle perimeter-over-six.

  \item[Ortho-unit polygons, point guards.]
  Proven tight: every ortho-unit polygon with $n \geq 12$ vertices is guardable
  by $\floor{(n-4)/8}$ guards \cite{DiazBanezEtAl2025}.  Since $n=P$ in the
  ortho-unit model, this is stronger than $P/6$ for that narrow class.

  \item[Integral orthogonal polygons, point guards, perimeter $N$.]
  Hole-free case proven via Ma{\ss}berg: after subdividing long boundary edges
  into unit grid edges, a hole-free integral orthogonal polygon is a hole-free
  polyomino with $N=\ell$.  Diaz-Banez et al. state the natural $N/6$
  conjecture in the broader integral-perimeter discussion
  \cite{DiazBanezEtAl2025}; its unresolved content is therefore beyond the
  hole-free case, notably holes.

  \item[Polycubes and polyhypercubes.]
  Pinciu proves the higher-dimensional cell-count analogue: every
  $m$-polyhypercube in dimension $d\geq 2$ is guardable by
  $\floor{(m+1)/3}$ point guards, and some examples require that many
  \cite{Pinciu2015}.  For $d=3$ this is the polycube version of the
  polyomino area theorem.  It is still a unit-cell theorem, not a theorem for
  arbitrary continuous orthogonal polyhedra.

  \item[Orthogonal polyhedra, point guards.]
  No planar-style linear theorem exists in terms of vertices.  The Paterson--Yao
  bound gives $\Theta(n^{3/2})$ point guards as sufficient and sometimes
  necessary for $n$-vertex orthogonal polyhedra \cite{PatersonYao1992,VigliettaThesis}.

  \item[Orthogonal polyhedra, edge and reflex-edge guards.]
  The central open direction is an Urrutia-type edge-guard theorem:
  lower-bound examples require roughly $m/12$ edge guards, while known upper
  bounds are larger.  Benbernou et al. prove that $\floor{(m+r)/12}$ open
  edge guards suffice for an orthogonal polyhedron with $m$ edges and $r$
  reflex edges; their corollaries give $(11/72)m-g/6-1$ and
  $(7/12)r-g+1$ open-edge-guard bounds
  \cite[Theorem~5 and Corollaries~6--7]{BenbernouEtAl2011}.  The latter is a
  bound parameterized by reflex-edge count, not a reflex-edge-only guard
  theorem.  Viglietta proves optimal-style bounds for
  2-reflex orthogonal polyhedra \cite{Viglietta2020}.  Cano, Toth, Urrutia,
  and Viglietta improve the general 3D polyhedron edge-guard bound to
  $(5/6)m$, while noting a substantial remaining gap
  \cite{CanoTothUrrutiaViglietta2022}.  Aldana-Galvan et al. prove a related
  $\pi/2$-edge-guard theorem for orthogonal polyhedra
  \cite{AldanaGalvanEtAl2016}; this is adjacent limited-field context, not a
  standard edge-guard or reflex-edge-guard closure.
\end{description}

\section{Two-Dimensional Orthogonal Polygons}

\subsection{Classical Background}

The original art-gallery theorem is due to Chvatal: $\floor{n/3}$ guards
always suffice for a simple $n$-vertex polygon, and this is tight
\cite{Chvatal1975}.  Fisk's short proof via triangulation and three-coloring
became the standard proof template \cite{Fisk1978}.  O'Rourke's book remains
the standard comprehensive reference for the classical theory
\cite{ORourke1987}; Shermer and Urrutia survey art-gallery and illumination
problems, including variants and open problems \cite{Shermer1992,Urrutia2000}.
On the algorithmic side, Abrahamsen, Adamaszek, and Miltzow prove that the
general art-gallery problem is $\exists\mathbb{R}$-complete even for simple
polygons with integer-coordinate vertices, underscoring why restricted
orthogonal and lattice models are usually treated separately
\cite[Theorem~1]{AbrahamsenAdamaszekMiltzow2022}.

\subsection{Simple Orthogonal Polygons}

Kahn, Klawe, and Kleitman prove the tight orthogonal theorem:
$\floor{n/4}$ point guards suffice for every simple orthogonal polygon with
$n$ vertices \cite{KahnKlaweKleitman1983}.  Gyori gives a short proof
\cite{Gyori1986}.  This theorem is a vertex-count theorem, not a perimeter
theorem.  In a lattice setting the vertex count and the perimeter may be very
different, which is why the polyomino and integral-perimeter questions are
separate.

The proof tradition connects to decomposition, coloring, and perfect-graph
methods.  Motwani, Raghunathan, and Saran study covering orthogonal polygons
with star polygons \cite{MotwaniRaghunathanSaran1990}; Worman and Keil develop
decomposition and visibility-graph machinery for orthogonal art-gallery
problems \cite{WormanKeil2007}.  Bj{\o}rling-Sachs proves the tight
$\floor{n/4}$ edge-guard bound for rectilinear polygons, a useful 2D reference
point for later edge-guard models \cite{BjorlingSachs1998}.

\subsection{Holes}

Holes are one of the main fault lines in the literature.  For orthogonal
polygons with $h$ holes, classical vertex-guard bounds exist, but the strongest
natural conjectures are not settled in general.  Urrutia's open-problem list
records Shermer's conjecture that an $n$-vertex orthogonal polygon with $h$
holes should be guardable by $\floor{(n+h)/4}$ vertex guards, and Hoffmann's
conjecture that $\floor{2n/7}$ vertex guards should suffice for orthogonal
polygons with holes \cite{UrrutiaOpenProblems}.  Zylinski proves the
$\floor{(n+h)/4}$ bound under a cactus-dual quadrilateralization condition and
recovers Aggarwal's theorem for $h \leq 2$ \cite{Zylinski2006}.  Hoffmann and
Kriegel prove graph-coloring consequences for rectilinear polygons with holes,
including upper bounds for rectilinear prison-yard variants
\cite{HoffmannKriegel1996}.  Michael and Pinciu later reframe much of this
orthogonal vertex-guard theory through same-sign diagonal graphs and vertex
covers; in particular, they give an h-independent upper bound
$\floor{(17n-8)/52}$ for orthogonal polygons with any number of holes
\cite{MichaelPinciu2016}.

These hole results are relevant context, but they do not prove the
perimeter-over-six statement.  They usually concern vertex guards and vertex
counts; the perimeter-over-six problem concerns point guards and lattice
perimeter.

\section{Polyomino and Perimeter Results}

\subsection{Hole-Free Polyominoes}

Ma{\ss}berg proves the direct perimeter theorem for hole-free polyominoes
\cite[Theorem~1]{Massberg2014}.  In the notation of that paper, if a
hole-free polyomino has lattice perimeter $\ell$, then it can be guarded by at
most $\max\{1,\loversix\}$ point guards; for $\ell \geq 6$ this is exactly
$\loversix$.  The proof uses maximal axis-parallel rectangles in a rectilinear
gallery, relates a rectangle packing number to the guard number, and then
bounds the packing by the perimeter.  A comb construction shows tightness
\cite[Fig.~3 and p.~575]{Massberg2014}.

\subsection{Polyominoes with Holes}

The area theorem does allow holes: Biedl, Irfan, Iwerks, Kim, and Mitchell
prove that every connected $m$-cell polyomino with $m\geq 2$ can be guarded
by $\floor{(m+1)/3}$ point guards, and the bound is tight
\cite[Cor.~1]{BiedlEtAl2012}.  The one-cell case is trivial, and the
sufficiency proof works even under restrictive $r$-visibility
\cite[Theorem~1]{BiedlEtAl2012}.  This is the strongest clean theorem for
arbitrary polyominoes with holes in the standard point-guard model, but it is
an area bound rather than a perimeter bound.

For perimeter, the hole case remains conditional.  Ma{\ss}berg observes that
the perimeter theorem would extend to holes if a maximal-rectangle packing
conjecture were true \cite[Conj.~6.10]{MassbergHabilitation}.  The thesis
also proves the perimeter charging lemma for maximum packings in polyominoes
with possible holes \cite[Lem.~6.9]{MassbergHabilitation}.  This is
\emph{not} a proof of perimeter-over-six for holes; it identifies the missing
structural statement.

\subsection{Ortho-Unit and Integral Orthogonal Polygons}

Diaz-Banez et al. prove that every ortho-unit polygon with $n \geq 12$ vertices
can be guarded by $\floor{(n-4)/8}$ guards, and that the bound is tight
\cite{DiazBanezEtAl2025}.  Since every edge has unit length in an ortho-unit
polygon, $n$ is the perimeter.  This solves a narrow unit-edge class more
strongly than $P/6$.

The same paper frames the integral-perimeter conjecture: every integral
orthogonal polygon of perimeter $N$ should be guardable by $\floor{N/6}$ point
guards.  If ``integral orthogonal polygon'' is read in the usual simple,
hole-free sense, this is already covered by Ma{\ss}berg.  Subdividing long
sides into unit grid edges turns the polygon into a hole-free polyomino with
the same lattice perimeter, $N=\ell$, and Ma{\ss}berg proves
$\max\{1,\floor{\ell/6}\}$ for that class
\cite[Theorem~1]{Massberg2014}.  Thus the Diaz-Banez question is open only
under a broader reading, such as integral domains with holes or another
lattice-domain convention not reduced to a hole-free polyomino.

\section{Three-Dimensional Orthogonal Polyhedra}

\subsection{Point Guards}

The direct 3D point-guard analogue of planar art-gallery theorems fails.  In
orthogonal polyhedra with $n$ vertices, Paterson and Yao's binary-space
partition result implies a tight $\Theta(n^{3/2})$ point-guard bound
\cite{PatersonYao1992}.  Viglietta's thesis presents this as Theorem 3.7 and
draws the key modeling conclusion: point guards do not yield a linear
orthogonal-polyhedron art-gallery theorem \cite{VigliettaThesis}.  Therefore a
3D ``perimeter-over-six'' theorem is not obtained by simply replacing polygons
with orthogonal polyhedra and keeping point guards.

\subsection{Edge and Reflex-Edge Guards}

The 3D literature often treats edge guards as the analogue of planar vertex
guards.  Urrutia conjectured that an orthogonal polyhedron with $m$ edges
should be guardable by $m/12+O(1)$ closed edge guards, matching known lower
bound examples up to an additive constant \cite{VigliettaThesis}.  A related
reflex-edge conjecture asks whether roughly $(r-g)/2+1$ reflex edge guards
suffice, where $r$ is the number of reflex edges and $g$ is the genus.

Benbernou, Demaine, Demaine, Kurdia, O'Rourke, Toussaint, Urrutia, and
Viglietta improve the known asymptotic upper bounds for orthogonal polyhedra:
their main open-edge theorem gives $\floor{(m+r)/12}$ guards, and their
corollaries yield $(11/72)m-g/6-1$ in terms of total edges and
$(7/12)r-g+1$ in terms of reflex-edge count
\cite[Theorem~5 and Corollaries~6--7]{BenbernouEtAl2011}.  This is an
edge-guard theorem parameterized by $r$, not a reflex-edge-only guard theorem.
Viglietta later proves
optimal-style bounds for 2-reflex orthogonal polyhedra, namely orthogonal
polyhedra whose reflex edges occur in only two directions
\cite{Viglietta2020}.  These are major partial results, but they do not settle
the general 3-reflex orthogonal-polyhedron conjectures.

Cano, Toth, Urrutia, and Viglietta give a broad theorem for arbitrary
polyhedra in three-space: every polyhedron with $m$ edges can be guarded by at
most $(5/6)m$ edge guards \cite{CanoTothUrrutiaViglietta2022}.  This improves
the general-polyhedron edge-guard bound, but it still leaves a large gap from
the conjectured $m/6$ style lower/upper behavior for general polyhedra and the
$m/12+O(1)$ orthogonal-polyhedron conjecture.

Aldana-Galvan, Alvarez-Rebollar, Catana-Salazar, Jimenez-Salinas,
Solis-Villarreal, and Urrutia prove a theorem for $\pi/2$-edge guards in
orthogonal polyhedra \cite{AldanaGalvanEtAl2016}.  This source belongs in the
3D edge-guard neighborhood, but the guard has limited field of view, so it is
recorded as a separate adjacent model rather than evidence for the standard
closed-edge or reflex-edge conjectures.

\subsection{Face Guards and Polycube Variants}

Face guards are another 3D model.  Viglietta studies face-guarding for several
classes of polyhedra, including orthogonal polyhedra, 4-oriented polyhedra, and
2-reflex orthostacks \cite{VigliettaFace2014}.  Face guards are useful for
mapping the 3D model space, but they are not a direct continuation of the
point-guard perimeter-over-six question.

For lattice 3D objects, Pinciu's theorem gives the direct point-guard
cell-volume analogue: $\floor{(m+1)/3}$ point guards are tight for
$m$-polyhypercubes in every dimension $d\geq 2$ \cite{Pinciu2015}.  Pinciu's
separate pixel/voxel-guard bounds remain a discrete cell-guard line with an
explicit sharpness conjecture, so the paper keeps point guards and cell guards
in separate model cells.

\section{Algorithmic and Optimization Context}

Extremal art-gallery theorems answer a worst-case existence question: how many
guards always suffice for every object in a class?  Optimization papers ask a
different question: given one object, how hard is it to find a minimum guard
set?  The distinction matters because a clean universal upper bound need not
yield an efficient exact algorithm for the minimum set, and an NP-hardness
result does not rule out a sharp extremal theorem.

Lee and Lin established the classical NP-hardness framework for art-gallery
optimization problems \cite{LeeLin1986}.  In orthogonal settings,
Schuchardt--Hecker prove NP-hardness for two orthogonal art-gallery problems
\cite{SchuchardtHecker1995}, and Katz--Roisman study hardness for guarding
vertices of rectilinear domains \cite{KatzRoisman2008}.  At the level of the
general point-guard problem, Abrahamsen--Adamaszek--Miltzow show
$\exists\mathbb{R}$-completeness, already for simple polygons with
integer-coordinate vertices \cite[Theorem~1]{AbrahamsenAdamaszekMiltzow2022}.
These complexity results should be read as optimization barriers, not as
barriers to universal counting theorems such as $\floor{n/4}$ or
$\floor{\ell/6}$.

Approximation and restricted-visibility models form a parallel literature.
Ghosh gives broad approximation-algorithm context for art-gallery problems
\cite{Ghosh2010}.  Durocher et al. and Biedl et al. study sliding cameras and
sliding $k$-transmitters in orthogonal polygons, including holes and hardness
or approximation results in those models
\cite{DurocherEtAl2017,BiedlEtAl2019}.  Filtser et al. study polyominoes under
$k$-hop visibility, a discrete variant close to network or city-block
applications rather than straight-line point-guard visibility
\cite{FiltserEtAl2025}.  These variants are useful sources of techniques and
open problems, but they should not be cited as settling the standard
perimeter-over-six question.

\section{Open Problems and Research Directions}

The following problems are the most directly connected to the survey above.
Items are grouped by mathematical model rather than by perceived difficulty.
When a question is only a consequence of comparing known results, the text
says so explicitly; otherwise the cited source records the conjecture or the
conditional dependency.

Two entries below are different in kind.  They are recorded as working prompts
for Galen from the current project, not as literature-confirmed open problems:
Paul's flat-vertex hybrid expression and a possible 3D surface-area-over-eight
analogue.  They are included because they are useful research directions, but
the table and prose mark the missing model choices explicitly.

\begin{table}[!t]
\small
\setlength{\tabcolsep}{3pt}
\renewcommand{\arraystretch}{1.12}
\begin{tabularx}{\textwidth}{
  >{\raggedright\arraybackslash}p{0.22\textwidth}
  >{\raggedright\arraybackslash}p{0.20\textwidth}
  Y
  Y}
\toprule
Problem cluster & Status source & What is known & What would constitute progress \\
\midrule
Polyomino perimeter with holes &
Ma{\ss}berg habilitation \cite[Conj.~6.10]{MassbergHabilitation} &
Hole-free theorem is published; hole extension is conditional on a
maximal-rectangle packing conjecture. &
Prove the packing conjecture, prove $\floor{\ell/6}$ by another method, or
construct a hole counterexample. \\

Integral perimeter beyond the hole-free case &
Diaz-Banez et al.; Ma{\ss}berg \cite{DiazBanezEtAl2025,Massberg2014} &
Hole-free integral polygons reduce to hole-free polyominoes. &
State and resolve the exact holes/broader-domain version, including how
perimeter is counted. \\

Paul/Galen flat-vertex hybrid &
Project working question, not source-stated &
If $P=n+f$, then $n/8+f/4<P/6$ is equivalent to $f<n/2$; no source was found
for a guard theorem using this weighted boundary 1-skeleton parameter. &
Fix $f$, holes, and guard model; decide whether the target is an upper bound,
lower-bound family, or improvement regime. \\

Orthogonal polygons with holes, vertex guards &
Urrutia; Zylinski; Michael--Pinciu
\cite{UrrutiaOpenProblems,Zylinski2006,MichaelPinciu2016} &
Strong conjectures remain open; partial and h-independent upper bounds are
known. &
Close the gap between lower-bound constructions and the best vertex-guard
upper bounds. \\

3D orthogonal edge guards &
Viglietta; Benbernou et al. \cite{VigliettaThesis,BenbernouEtAl2011} &
Known upper bounds exceed Urrutia-type lower-bound targets. &
Prove an $m/12+O(1)$ style theorem or identify a larger necessary correction
term. \\

3D reflex-edge guards &
Viglietta \cite{Viglietta2020} &
2-reflex orthogonal polyhedra are understood much better than general
3-reflex orthogonal polyhedra. &
Extend the induction/charging methods beyond the 2-reflex restriction. \\

3D surface-area-over-eight &
Project working conjecture, not source-stated &
The question needs lattice-normalized surface area; in a unit boundary mesh
$S/8=E_{\partial}/16$, but unrefined edge count is a different parameter. &
Choose object class and guard model, then test lower- and upper-bound
constructions. \\

Polycube/polyhypercube voxel guards &
Pinciu \cite{Pinciu2015} &
Point-guard cell-volume bounds are tight, but pixel/voxel guard bounds have
only partial formulas and a sharpness conjecture. &
Prove or refute the dependent-on-$d$ pixel/voxel-guard conjecture. \\

Algorithmic restricted cases &
Lee--Lin; Schuchardt--Hecker; Ghosh
\cite{LeeLin1986,SchuchardtHecker1995,Ghosh2010} &
Exact optimization is hard in many settings; approximation and special
visibility models remain active. &
Find tractable subclasses aligned with the extremal perimeter or edge-guard
conjectures. \\
\bottomrule
\end{tabularx}
\caption{Research-facing status table for the main open-problem clusters.}
\label{tab:open-problem-map}
\end{table}

\begin{enumerate}
  \item \textbf{Perimeter-over-six for polyominoes with holes.}
  Let $P$ be a polyomino with holes and total lattice perimeter $\ell$,
  counting both outer and hole boundaries.  Is $P$ always guardable by
  $\max\{1,\loversix\}$ point guards?  Ma{\ss}berg's habilitation gives a
  conditional route to this theorem through a maximal-rectangle packing
  conjecture \cite[Lem.~6.9 and Conj.~6.10]{MassbergHabilitation}; the
  published perimeter theorem is the hole-free case
  \cite[Theorem~1]{Massberg2014}.

  \item \textbf{Ma{\ss}berg's maximal-rectangle packing conjecture.}
  In a rectilinear gallery with holes, does the maximum packing size of
  maximal rectangles control the required number of point guards?  This is the
  structural conjecture identified by Ma{\ss}berg as sufficient for extending
  the perimeter-over-six theorem to polyominoes with holes
  \cite[Conj.~6.10]{MassbergHabilitation}.  The same discussion notes that
  maximal-rectangle intersection graphs need not be perfect in the hole case,
  so the conjecture is not a routine extension of the hole-free perfect-graph
  argument.

  \item \textbf{Integral-perimeter bounds beyond the hole-free case.}
  The hole-free integral orthogonal polygon version follows from
  Ma{\ss}berg by unit-grid subdivision.  The remaining integral-perimeter
  question is whether the same $\floor{N/6}$ behavior survives when holes or
  broader lattice-domain conventions are allowed
  \cite{Massberg2014,DiazBanezEtAl2025}.

  \item \textbf{Paul/Galen flat-vertex hybrid question.}
  This is a project working prompt rather than a source-stated open problem.
  Suppose an integral orthogonal polygon has been subdivided into unit boundary
  steps, with total lattice perimeter $P$, $n$ turn vertices, and $f$ flat
  subdivision vertices.  Under the convention $P=n+f$,
  \[
    n/8+f/4<P/6 \quad\Longleftrightarrow\quad f<n/2.
  \]
  Thus the inequality is not itself the mathematical content.  The real
  question is whether the hybrid expression $n/8+f/4$ controls a guard
  theorem, lower-bound family, or useful improvement regime once $f$, holes,
  and guard model are fixed.  In 1-skeleton language, the unit-subdivided
  boundary cycle has $P$ unit edges and $P$ boundary vertices; the vertices
  are partitioned into $n$ turn vertices and $f$ flat subdivision vertices.
  So this prompt asks whether weighting a single boundary skeleton differently
  at turn and straight-through vertices has guard-theoretic content.  The
  current evidence pass found no primary source for an $n/8+f/4$ art-gallery
  theorem or source-stated open problem.

  \item \textbf{Vertex guards for orthogonal polygons with holes.}
  Shermer's conjecture asks whether every $n$-vertex orthogonal polygon with
  $h$ holes is guardable by $\floor{(n+h)/4}$ vertex guards.  Hoffmann's
  conjecture asks whether $\floor{2n/7}$ vertex guards always suffice.  These
  are source-stated open problems in the orthogonal-with-holes literature, and
  they are related to but distinct from perimeter-type point-guard bounds
  \cite{UrrutiaOpenProblems,HoffmannKriegel1996,Zylinski2006,MichaelPinciu2016}.

  \item \textbf{A correction term if perimeter-over-six fails with holes.}
  This is a derived direction rather than a named conjecture.  If holes force
  counterexamples to a pure $\ell/6$ theorem, a useful replacement may need to
  account for total perimeter, outer perimeter, number of holes, hole
  perimeter, reflex vertices, or rectangle-packing obstructions.

  \item \textbf{Urrutia-type edge-guard bounds for orthogonal polyhedra.}
  For genus-zero orthogonal polyhedra with $m$ edges, the central conjectural
  target is an $m/12+O(1)$ edge-guard bound.  Lower-bound examples and current
  upper bounds leave a substantial gap \cite{VigliettaThesis,BenbernouEtAl2011}.

  \item \textbf{Reflex-edge guards in general orthogonal polyhedra.}
  Viglietta proves tight-style bounds for 2-reflex orthogonal polyhedra, where
  reflex edges occur in only two coordinate directions, but the analogous
  reflex-edge conjecture for general orthogonal polyhedra remains open
  \cite{Viglietta2020}.

  \item \textbf{The general three-dimensional edge-guard gap.}
  Cano, Toth, Urrutia, and Viglietta improve the edge-guard upper bound for
  arbitrary polyhedra to $(5/6)m$, but explicitly leave the gap between known
  lower and upper bounds as an open problem
  \cite{CanoTothUrrutiaViglietta2022}.  Orthogonal polyhedra form the more
  structured subclass where sharper Urrutia-type conjectures are expected.

  \item \textbf{Surface-area-over-eight in 3D.}
  This is also a project working conjecture rather than a literature-confirmed
  theorem.  A meaningful statement must first choose a lattice-normalized
  surface area $S$, such as the number of exposed unit square faces of a
  polycube or grid-refined orthogonal polyhedron, and must specify whether the
  guards are points, vertices, edges, faces, or unit-face/cell guards.  Nearby
  sources do not settle this: Viglietta's face-guard paper uses a different
  guard model \cite{VigliettaFace2014}, Pinciu's polyhypercube theorem is a
  volume/cell-count point-guard result \cite{Pinciu2015}, and continuous 3D
  point guarding has different vertex-count behavior
  \cite{PatersonYao1992,VigliettaThesis}.

  The same prompt can be phrased through the unit-refined boundary
  1-skeleton.  Let $E_{\partial}$ be the number of unit edges in the boundary
  surface mesh.  For a closed unit-square boundary surface, each square
  contributes four edge incidences and each boundary edge is incident to two
  squares, so $E_{\partial}=2S$.  In that refined model the target $S/8$ is
  exactly $E_{\partial}/16$.  This is not the unrefined edge count $m$ used in
  continuous edge-guard results: scaling an orthogonal box changes Euclidean
  surface area while leaving the unrefined polyhedron graph unchanged.  If
  ``size of the 1D skeleton'' means $V_{\partial}+E_{\partial}$ instead of
  edge count alone, the connected orientable case also carries the Euler term
  $V_{\partial}+E_{\partial}=3S+2-2g$, so the convention should be stated
  before this becomes a theorem-style conjecture.

  \item \textbf{Pixel/voxel guards for polyhypercubes.}
  Pinciu settles the point-guard cell-volume question for every dimension
  $d\geq 2$, but the corresponding pixel/voxel-guard model is different.  In
  that model Pinciu gives lower bounds, dimension-independent upper bounds, and
  a conjectured sharp dependent-on-$d$ formula \cite{Pinciu2015}.  This is an
  open problem for discrete cell guards, not a gap in the point-guard
  $\floor{(m+1)/3}$ theorem.

  \item \textbf{Algorithmic complexity and approximation in restricted models.}
  Extremal guard bounds and optimization algorithms remain separate questions.
  Minimum guarding is NP-hard in many orthogonal settings
  \cite{LeeLin1986,SchuchardtHecker1995,KatzRoisman2008,BiedlEtAl2012,DurocherEtAl2017,BiedlEtAl2019}.
  A parallel algorithmic direction is to identify exact polynomial subclasses
  and approximation guarantees for polyominoes with holes, integral orthogonal
  domains, polycubes, and orthogonal polyhedra.  Recent work on $k$-hop
  visibility in polyominoes records related approximation questions
  \cite{FiltserEtAl2025}.
\end{enumerate}

\section{Annotated Source Guide}

\subsection*{Foundational Art-Gallery Theory}

\begin{description}
  \item[Chvatal \cite{Chvatal1975} and Fisk \cite{Fisk1978}.]
  Foundational $\floor{n/3}$ theorem for simple polygons and the short
  triangulation/coloring proof.

  \item[O'Rourke \cite{ORourke1987}.]
  Standard book reference for art-gallery theorems, algorithms, visibility,
  decompositions, and early 3D context.

  \item[Shermer \cite{Shermer1992} and Urrutia \cite{Urrutia2000}.]
  Survey references for art-gallery and illumination problems, useful for
  terminology, variants, and classical open-problem context.

  \item[Lee--Lin \cite{LeeLin1986}.]
  Classical computational-complexity source for art-gallery optimization
  problems.

  \item[Abrahamsen--Adamaszek--Miltzow \cite{AbrahamsenAdamaszekMiltzow2022}.]
  Modern complexity reference proving $\exists\mathbb{R}$-completeness for the
  general art-gallery problem, even for simple integer-coordinate polygons.
\end{description}

\subsection*{2D Orthogonal Polygons and Holes}

\begin{description}
  \item[Kahn--Klawe--Kleitman \cite{KahnKlaweKleitman1983}.]
  Primary source for the tight $\floor{n/4}$ theorem for simple orthogonal
  polygons.

  \item[Gyori \cite{Gyori1986}.]
  Short proof of the rectilinear art-gallery theorem.

  \item[Hoffmann--Kriegel \cite{HoffmannKriegel1996}.]
  Graph-coloring result with consequences for rectilinear polygons with holes
  and prison-yard variants.

  \item[Zylinski \cite{Zylinski2006}.]
  Coloring proof of Aggarwal's theorem for orthogonal galleries with holes,
  including the $\floor{(n+h)/4}$ bound for $h\leq 2$ and cactus-dual cases.

  \item[Michael--Pinciu \cite{MichaelPinciu2016}.]
  Same-sign diagonal graph and vertex-cover framework for orthogonal polygons,
  including h-independent bounds for polygons with holes.

  \item[Urrutia open-problem page \cite{UrrutiaOpenProblems}.]
  Direct source for Shermer's and Hoffmann's open conjectures on orthogonal
  polygons with holes using vertex guards.

  \item[Schuchardt--Hecker \cite{SchuchardtHecker1995}.]
  NP-hardness source for art-gallery problems in orthogonal polygons.

  \item[Worman--Keil \cite{WormanKeil2007}.]
  Decomposition and visibility-graph machinery for the orthogonal art-gallery
  problem.

  \item[Bj{\o}rling-Sachs \cite{BjorlingSachs1998}.]
  Tight edge-guard theorem for rectilinear polygons; relevant when comparing
  2D orthogonal edge guards with 3D edge-guard conjectures.

  \item[Katz--Roisman \cite{KatzRoisman2008}.]
  NP-hardness for guarding vertices of rectilinear domains by vertex, boundary,
  or point guards.
\end{description}

\subsection*{Polyomino, Perimeter, and Lattice Variants}

\begin{description}
  \item[Ma{\ss}berg \cite{Massberg2014}.]
  Primary peer-reviewed source for perimeter-over-six in hole-free
  polyominoes: Theorem 1 gives $\max\{1,\floor{\ell/6}\}$ and Fig. 3 gives the
  tight comb construction.  This also proves the nontrivial hole-free integral
  orthogonal polygon $\floor{N/6}$ statement after unit-grid subdivision.

  \item[Ma{\ss}berg habilitation \cite{MassbergHabilitation}.]
  Source for the conditional hole extension.  Lemma 6.9 gives the perimeter
  charging statement for maximum rectangle packings in polyominoes with
  possible holes, and Conjecture 6.10 is the missing guard-number statement
  for rectilinear galleries with holes.

  \item[Biedl--Irfan--Iwerks--Kim--Mitchell \cite{BiedlEtAl2012}.]
  Tight $\floor{(m+1)/3}$ cell-count theorem for connected $m$-polyominoes
  with $m\geq 2$, holes allowed.  The definitions are on p.~712, the comb
  lower bound is Fig.~3, and Corollary~1 gives the theorem.

  \item[Diaz-Banez et al. \cite{DiazBanezEtAl2025}.]
  Recent tight theorem for ortho-unit polygons and source of the integral
  orthogonal polygon $\floor{N/6}$ conjecture.  Its simple/hole-free reading
  overlaps with Ma{\ss}berg rather than creating a separate open problem.

  \item[Pinciu \cite{Pinciu2015}.]
  Point-guard cell-count theorem: for every $d\geq 2$,
  $\floor{(m+1)/3}$ guards are tight for $m$-poly\-hypercubes.  The same
  source also records separate pixel/voxel-guard bounds
  and a sharpness conjecture.

  \item[Filtser--Krohn--Nilsson--Rieck--Schmidt \cite{FiltserEtAl2025}.]
  Current work on $k$-hop visibility in polyominoes, including open
  approximation questions for polyominoes with holes.
\end{description}

\subsection*{3D Orthogonal Polyhedra}

\begin{description}
  \item[Paterson--Yao \cite{PatersonYao1992}.]
  Binary-space partition result used for the $\Theta(n^{3/2})$ point-guard
  bound in orthogonal polyhedra.

  \item[Viglietta thesis \cite{VigliettaThesis}.]
  Central survey and source for guarding/searching polyhedra, the point-guard
  obstruction, and Urrutia-type edge-guard conjectures.

  \item[Benbernou et al. \cite{BenbernouEtAl2011}.]
  CCCG source for improved edge-guard bounds in orthogonal polyhedra:
  $\floor{(m+r)/12}$ open edge guards, with corollaries
  $(11/72)m-g/6-1$ and $(7/12)r-g+1$.  The $r$ bound is parameterized by
  reflex-edge count; it is not a reflex-edge-only guard theorem.

  \item[Aldana-Galvan et al. \cite{AldanaGalvanEtAl2016}.]
  Adjacent CCCG source for $\pi/2$-edge guards in orthogonal polyhedra.  It is
  useful limited-field context, but it should not be read as a standard
  edge-guard or reflex-edge-guard theorem.

  \item[Viglietta 2-reflex paper \cite{Viglietta2020}.]
  Proves optimal-style reflex-edge bounds for 2-reflex orthogonal polyhedra and
  gives an $O(n\log n)$ guard-placement algorithm.

  \item[Viglietta face-guarding paper \cite{VigliettaFace2014}.]
  Face-guarding results for several classes of polyhedra, including orthogonal
  subclasses.

  \item[Cano--Toth--Urrutia--Viglietta \cite{CanoTothUrrutiaViglietta2022}.]
  General edge-guard result in three-space and explicit open gap.
\end{description}

\subsection*{Algorithms and Related Visibility Models}

\begin{description}
  \item[Motwani--Raghunathan--Saran \cite{MotwaniRaghunathanSaran1990}.]
  Perfect-graph approach to covering orthogonal polygons with star polygons.

  \item[Ghosh \cite{Ghosh2010}.]
  General approximation-algorithm context for art-gallery problems.

  \item[Durocher et al. \cite{DurocherEtAl2017}.]
  Sliding-camera model for orthogonal art galleries, including holes.

  \item[Biedl et al. \cite{BiedlEtAl2019}.]
  Sliding $k$-transmitter hardness and approximation results for orthogonal
  polygons.
\end{description}

\section{Completeness Protocol for This Review}

This review treats completeness as an auditable process rather than as a claim
that any finite search has exhausted the literature.  A result should enter
the narrative only after its source, exact model, theorem or problem statement,
safe translation, coverage status, and open-problem consequences have been
recorded.  This discipline matters here because small changes in holes,
lattice assumptions, guard type, dimension, and parameter frequently turn a
valid theorem into a non-comparable statement.

\begin{center}
\footnotesize
\renewcommand{\arraystretch}{1.08}
\begin{tabularx}{\textwidth}{
  >{\raggedright\arraybackslash}p{0.08\textwidth}
  >{\raggedright\arraybackslash}p{0.34\textwidth}
  Y}
\toprule
Step & Question & Review artifact \\
\midrule
1 & What is the exact scope cell?  Record object class, holes, guard model,
dimension, and parameter before comparing results. &
Coverage matrix row or new coverage cell. \\

2 & What are the seed sources?  Start from books, surveys, theses,
open-problem pages, and named theorem papers. &
Source-registry records with primary DOI, publisher, arXiv, or author links
when available. \\

3 & What searches have been run?  Use exact-title searches, DOI/arXiv
lookups, author pages, backward references, forward citation searches, venue
scans, and terminology variants. &
Search-log entries with included sources, exclusions, and follow-up actions. \\

4 & Is the source primary for the claim?  Prefer the paper that proves or
states the result; use surveys and books to orient the search and explain
history. &
Source type and role recorded in the source registry. \\

5 & What exactly is the claim?  Extract theorem, conjecture, lower bound,
hardness, algorithmic, survey, or open-problem records without changing the
model. &
Claim-registry entry with status, model, parameter, locator, and citation. \\

6 & What does the claim not imply?  Record blocked translations across holes,
polyominoes, integral polygons, guard models, and 3D analogues. &
Translation note in the claim registry and, when central, prose in the
model-separation section. \\

7 & Which status changed?  Update every affected coverage cell as proved,
tight, conditional, open, synthesized, out of scope, or still needing
follow-up. &
Coverage matrix with supporting claim identifiers and next action. \\

8 & Is an open problem source-stated, conditional, or derived?  Keep these
categories separate and record dependencies. &
Open-problem ledger linked to supporting sources and claims. \\

9 & Are locators and quotations sufficient?  Add theorem numbers, page
numbers, figure references, or short quotations for problem statements where
available. &
Evidence-report gap list should shrink after each locator pass. \\

10 & Does the document still assemble cleanly?  Check that the short answer,
status maps, proof discussion, open problems, source guide, and evidence trail
still form one reader-facing flow. &
Document assembly map and paper outline. \\

11 & Is the framework progress current?  Record which stage changed, what
remains incomplete, and which artifact should be updated next. &
Progress tracker and current work queue. \\

12 & Has a final status pass been run for central open, conditional, or
recent-status claims?  Repeat forward-citation, venue, author-page,
exact-title, and adjacent-model searches before preserving a status. &
Search-log and search-query entries with dated outcomes for each central
status cell. \\

13 & Is the prose synchronized with the evidence?  Regenerate the evidence
report, run registry checks, rebuild the PDF, and inspect the log before
finalizing. &
Updated paper, generated evidence report, and clean validation run. \\
\bottomrule
\end{tabularx}
\end{center}

\subsection*{Completeness Gates}

The practical standard for this paper is not ``no missing papers.''  It is
that every central status statement can be audited.  In particular, a claim
that a problem is open should be supported by one of three paths: an explicit
open-problem statement in a primary or survey source, a conditional theorem
whose stated hypothesis remains unproved in the tracked sources, or a derived
gap visible from the coverage matrix after the relevant model cells have been
searched.  The last category should be phrased cautiously and should identify
the completed searches that make the gap credible.

The current framework therefore makes gaps visible instead of hiding them in
prose.  At the time of this pass, every claim record has a locator at the
level currently available from primary PDFs, open-problem pages, public
technical reports, or primary metadata.  The forward/adjacent sweep and final
OpenAlex/web pass for the central holes, perimeter, and 3D edge/reflex-edge
clusters are now logged.  They did not add a tracked theorem that changes the
main statuses.  It also records two project working questions without
promoting them to literature-confirmed open problems: Paul's
$n/8+f/4<P/6$ flat-vertex prompt and a 3D surface-area-over-eight prompt.
The current coverage matrix now separates the cited adjacent variants that
already have source records: two-dimensional edge guards, sliding-camera and
sliding-transmitter visibility, $k$-hop visibility, polyhypercube point
guards, polyhypercube pixel/voxel guards, 3D face guards, $\pi/2$-edge
guards, and 2-reflex reflex-edge guards.  It also records scope-exclusion
cells for floodlight, half-plane,
dispersive, contiguous, mobile,
point-boundary, and related variants.  Those cells are not source claims; they
are checkpoints requiring source intake before any such variant is used as a
theorem-level result.  The remaining completeness work is recurring
search-frontier work and model-confirmation work rather than a missing-locator
problem.

\section{Evidence Audit Trail}

The review is backed by a structured evidence layer in the repository.  The
source registry records bibliographic information and links; the claims
registry records the specific theorem, conjecture, hardness, survey, or
open-problem statement extracted from each source; the coverage matrix records
which model cells are proved, open, conditional, synthesized, or still in need
of follow-up; and the open-problems ledger records the unresolved questions and
their supporting sources.

The generated evidence report is intended as an audit companion to the prose.
It lists central proved results, open and conditional problems, coverage cells,
missing page or theorem locators, and next actions for improving the review.
The document assembly map records how the evidence artifacts become a readable
paper, and the framework progress tracker records which part of the review
process is current, which part needs work, and which artifact should be touched
next.  The prose above should be updated only after the corresponding source,
claim, coverage, open-problem, assembly, and progress records support the
statement being added.

The May 23, 2026 search-log entries are the current checkpoint for
recent-status claims.  They record exact-title, OpenAlex forward-citation,
arXiv/DROPS/topic, venue, and adjacent-model queries for the
orthogonal-polygons-with-holes, polyomino-perimeter-with-holes, integral
perimeter, and 3D orthogonal-polyhedron clusters.  The result was mainly
status preservation: the central open and conditional cells remain open or
conditional in the tracked evidence.  The one added theorem source is
Aldana-Galvan et al.'s $\pi/2$-edge-guard result, which is recorded as a
separate adjacent model rather than as a closure of the standard 3D
edge/reflex-edge gaps.

The May 26, 2026 Galen/Paul working-question pass records two proposed
questions separately from source-stated open problems.  No source was found for
an $n/8+f/4$ guard theorem, an $n/8+f/4<P/6$ open problem, or a
surface-area-over-eight theorem for 3D orthogonal objects.  The evidence layer
therefore treats both prompts as model-definition tasks first: fix the
flat-vertex parameter and intended guard statement in 2D, and fix the surface
normalization, object class, guard model, and 1D-skeleton convention in 3D.

\begin{center}
\footnotesize
\renewcommand{\arraystretch}{1.08}
\begin{tabularx}{\textwidth}{
  >{\raggedright\arraybackslash}p{0.25\textwidth}
  >{\raggedright\arraybackslash}p{0.18\textwidth}
  Y}
\toprule
Cell checked & Outcome & Reader-facing consequence \\
\midrule
Polyominoes with holes, point guards, lattice perimeter $\ell$ &
Status preserved &
No tracked source proves the $\ell/6$ holes extension; it remains tied to
Ma{\ss}berg's conditional route. \\

Integral orthogonal polygons, point guards, perimeter $N$ &
Status preserved &
The hole-free reading remains covered by Ma{\ss}berg; broader readings with
holes remain open or conditional in the tracked evidence. \\

Orthogonal polygons with holes, vertex guards &
Status preserved &
Floodlight and half-plane hits remain adjacent exclusions; no standard
vertex-guard resolution was tracked. \\

Orthogonal polyhedra, standard edge/reflex-edge guards &
Adjacent source added; status preserved &
Aldana-Galvan et al. add $\pi/2$-edge-guard context, but the standard
edge/reflex-edge gaps remain open. \\

Paul/Galen flat-vertex and surface-area prompts &
Working questions recorded; no source-status change &
The prompts are useful research directions, but the PDF treats them as
model-definition tasks until the missing conventions or sources are supplied;
the skeleton notes distinguish unit-refined boundary graphs from unrefined
polyhedron edge counts.
\\
\bottomrule
\end{tabularx}
\end{center}

\section{Conclusion}

\begin{enumerate}
  \item The theorem ``$\ell/6$ guards suffice'' is established for
  \emph{hole-free polyominoes}.  The polyomino-with-holes version is
  conditional on Ma{\ss}berg's rectangle-packing conjecture in the sources
  surveyed.

  \item The statement ``an integral orthogonal polygon of perimeter $N$ is
  guardable by $\floor{N/6}$ guards'' has two different readings.  The
  simple/hole-free reading is already proved by Ma{\ss}berg, because it is
  exactly the hole-free polyomino case after unit-grid subdivision.  The holes
  or broader-domain reading remains open/conditional.

  \item The right 3D terminology is \emph{orthogonal polyhedron} or
  \emph{orthogonal polytope}.  The natural 3D guard models are not obtained by
  copying the 2D point-guard setting.

  \item A 3D point-guard analogue of the planar perimeter theorem is blocked by
  the $\Theta(n^{3/2})$ point-guard bound.  The central 3D problems instead
  concern edge guards, reflex-edge guards, face guards, or discrete polycube
  visibility models.

  \item The most direct open directions are the polyomino-with-holes
  perimeter bound, Ma{\ss}berg's rectangle-packing conjecture, the
  holes/broader-domain $\floor{N/6}$ extension, and the Urrutia-type
  $m/12+O(1)$ and reflex-edge conjectures for general orthogonal polyhedra.

  \item Paul's $n/8+f/4<P/6$ prompt and the 3D surface-area-over-eight prompt
  are now recorded as proposed working questions for Galen.  They should not be
  cited as literature-confirmed open problems until their model conventions or
  source status are fixed.  The skeleton formulation is included as a useful
  normalization: in a unit-square 3D boundary mesh, $S/8$ is the same target as
  $E_{\partial}/16$, but that refined skeleton count is not the unrefined edge
  count used in standard 3D edge-guard results.
\end{enumerate}

\begin{thebibliography}{99}

\bibitem{AbrahamsenAdamaszekMiltzow2022}
M. Abrahamsen, A. Adamaszek, and T. Miltzow.
\newblock The art gallery problem is $\exists\mathbb{R}$-complete.
\newblock \emph{Journal of the ACM}, 69(1), Article 4, 2022.
\newblock \url{https://doi.org/10.1145/3486220}

\bibitem{AldanaGalvanEtAl2016}
I. Aldana-Galvan, J. L. Alvarez-Rebollar, J. C. Catana-Salazar,
M. Jimenez-Salinas, E. Solis-Villarreal, and J. Urrutia.
\newblock Minimizing the solid angle sum of orthogonal polyhedra and guarding
them with $\pi/2$-edge guards.
\newblock In \emph{Proceedings of the 28th Canadian Conference on
Computational Geometry}, 2016.
\newblock \url{https://mazayjimenez.org/papers/minimizing_cccg2016.pdf}
\newblock \url{https://www.cccg.ca/proceedings/2016/proceedings2016.pdf}

\bibitem{BenbernouEtAl2011}
N. M. Benbernou, E. D. Demaine, M. L. Demaine, A. Kurdia,
J. O'Rourke, G. T. Toussaint, J. Urrutia, and G. Viglietta.
\newblock Edge-guarding orthogonal polyhedra.
\newblock In \emph{Proceedings of the 23rd Canadian Conference on Computational Geometry}, 461--466, 2011.
\newblock \url{https://erikdemaine.org/papers/EdgeGuarding_CCCG2011/}
\newblock \url{https://www.science.smith.edu/~jorourke/Papers/EdgeGuarding.pdf}

\bibitem{BiedlEtAl2012}
T. C. Biedl, M. T. Irfan, J. Iwerks, J. Kim, and J. S. B. Mitchell.
\newblock The art gallery theorem for polyominoes.
\newblock \emph{Discrete \& Computational Geometry}, 48, 711--720, 2012.
\newblock \url{https://doi.org/10.1007/s00454-012-9429-1}
\newblock \url{https://tildesites.bowdoin.edu/~mirfan/papers/Art_Gallery_DCG_2012.pdf}
\newblock \url{https://researchconnect.suny.edu/en/publications/the-art-gallery-theorem-for-polyominoes/}

\bibitem{BiedlEtAl2019}
T. Biedl, T. M. Chan, S. Lee, S. Mehrabi, F. Montecchiani, H. Vosoughpour, and Z. Yu.
\newblock Guarding orthogonal art galleries with sliding $k$-transmitters: hardness and approximation.
\newblock \emph{Algorithmica}, 81(1), 69--97, 2019.
\newblock \url{https://doi.org/10.1007/s00453-018-0433-6}

\bibitem{BjorlingSachs1998}
I. Bj{\o}rling-Sachs.
\newblock Edge guards in rectilinear polygons.
\newblock \emph{Computational Geometry: Theory and Applications}, 11(2), 111--123, 1998.
\newblock \url{https://doi.org/10.1016/S0925-7721(98)00024-8}

\bibitem{CanoTothUrrutiaViglietta2022}
J. Cano, C. D. Toth, J. Urrutia, and G. Viglietta.
\newblock Edge guards for polyhedra in three-space.
\newblock \emph{Computational Geometry: Theory and Applications}, 104, 101859, 2022.
\newblock \url{https://doi.org/10.1016/j.comgeo.2022.101859}
\newblock \url{https://giovanniviglietta.com/papers/edgeguards.pdf}

\bibitem{Chvatal1975}
V. Chvatal.
\newblock A combinatorial theorem in plane geometry.
\newblock \emph{Journal of Combinatorial Theory, Series B}, 18(1), 39--41, 1975.
\newblock \url{https://doi.org/10.1016/0095-8956(75)90061-1}

\bibitem{DiazBanezEtAl2025}
J. M. Diaz-Banez, P. Horn, M. A. Lopez, N. Marin, A. Ramirez-Vigueras,
O. Sole-Pi, A. Stevens, and J. Urrutia.
\newblock Ortho-unit polygons can be guarded with at most $\floor{(n-4)/8}$ guards.
\newblock \emph{Graphs and Combinatorics}, 41, 15, 2025.
\newblock \url{https://doi.org/10.1007/s00373-024-02880-8}
\newblock \url{https://arxiv.org/abs/2208.12864}

\bibitem{DurocherEtAl2017}
S. Durocher, O. Filtser, R. Fraser, A. D. Mehrabi, and S. Mehrabi.
\newblock Guarding orthogonal art galleries with sliding cameras.
\newblock \emph{Computational Geometry: Theory and Applications}, 65, 12--26, 2017.
\newblock \url{https://doi.org/10.1016/j.comgeo.2017.04.001}

\bibitem{FiltserEtAl2025}
O. Filtser, E. Krohn, B. J. Nilsson, C. Rieck, and C. Schmidt.
\newblock Guarding polyominoes under $k$-hop visibility.
\newblock \emph{Algorithmica}, 87, 572--593, 2025.
\newblock \url{https://doi.org/10.1007/s00453-024-01292-7}

\bibitem{Fisk1978}
S. Fisk.
\newblock A short proof of Chvatal's watchman theorem.
\newblock \emph{Journal of Combinatorial Theory, Series B}, 24(3), 374, 1978.
\newblock \url{https://doi.org/10.1016/0095-8956(78)90059-X}

\bibitem{Ghosh2010}
S. K. Ghosh.
\newblock Approximation algorithms for art gallery problems in polygons.
\newblock \emph{Discrete Applied Mathematics}, 158(6):718--722, 2010.
\newblock \url{https://doi.org/10.1016/j.dam.2009.12.004}

\bibitem{Gyori1986}
E. Gyori.
\newblock A short proof of the rectilinear art gallery theorem.
\newblock \emph{SIAM Journal on Algebraic and Discrete Methods}, 7(3), 452--454, 1986.
\newblock \url{https://doi.org/10.1137/0607051}

\bibitem{HoffmannKriegel1996}
F. Hoffmann and K. Kriegel.
\newblock A graph-coloring result and its consequences for polygon-guarding problems.
\newblock \emph{SIAM Journal on Discrete Mathematics}, 9(2), 210--224, 1996.
\newblock \url{https://doi.org/10.1137/S0895480194265611}
\newblock \url{https://refubium.fu-berlin.de/handle/fub188/18382}

\bibitem{KahnKlaweKleitman1983}
J. Kahn, M. Klawe, and D. Kleitman.
\newblock Traditional galleries require fewer watchmen.
\newblock \emph{SIAM Journal on Algebraic and Discrete Methods}, 4(2), 194--206, 1983.
\newblock \url{https://doi.org/10.1137/0604020}

\bibitem{KatzRoisman2008}
M. J. Katz and G. S. Roisman.
\newblock On guarding the vertices of rectilinear domains.
\newblock \emph{Computational Geometry: Theory and Applications}, 39(3), 219--228, 2008.
\newblock \url{https://doi.org/10.1016/j.comgeo.2007.02.002}

\bibitem{LeeLin1986}
D. T. Lee and A. K. Lin.
\newblock Computational complexity of art gallery problems.
\newblock \emph{IEEE Transactions on Information Theory}, 32(2), 276--282, 1986.
\newblock \url{https://doi.org/10.1109/TIT.1986.1057165}

\bibitem{Massberg2014}
J. Ma{\ss}berg.
\newblock Perfect graphs and guarding rectilinear art galleries.
\newblock \emph{Discrete \& Computational Geometry}, 51, 569--577, 2014.
\newblock \url{https://doi.org/10.1007/s00454-014-9587-4}

\bibitem{MassbergHabilitation}
J. Ma{\ss}berg.
\newblock \emph{Geometrical and Combinatorial Optimization Problems}.
\newblock Habilitationsschrift, Universit\"at Ulm, 2015.
\newblock \url{https://www.uni-ulm.de/fileadmin/website_uni_ulm/mawi.inst.080/Massberg/habilitation_massberg.pdf}

\bibitem{MichaelPinciu2016}
T. S. Michael and V. Pinciu.
\newblock How to guard orthogonal polygons: diagonal graphs and vertex covers.
\newblock \emph{Discrete \& Computational Geometry}, 55(2), 410--422, 2016.
\newblock \url{https://doi.org/10.1007/s00454-015-9756-0}

\bibitem{MotwaniRaghunathanSaran1990}
R. Motwani, A. Raghunathan, and H. Saran.
\newblock Covering orthogonal polygons with star polygons: The perfect graph approach.
\newblock \emph{Journal of Computer and System Sciences}, 40(1), 19--48, 1990.
\newblock \url{https://doi.org/10.1016/0022-0000(90)90017-F}

\bibitem{ORourke1987}
J. O'Rourke.
\newblock \emph{Art Gallery Theorems and Algorithms}.
\newblock Oxford University Press, 1987.
\newblock \url{https://www.science.smith.edu/~jorourke/books/ArtGalleryTheorems/Art_Gallery_Full_Book.pdf}

\bibitem{PatersonYao1992}
M. S. Paterson and F. F. Yao.
\newblock Optimal binary space partitions for orthogonal objects.
\newblock \emph{Journal of Algorithms}, 13(1), 99--113, 1992.
\newblock \url{https://doi.org/10.1016/0196-6774(92)90007-Y}

\bibitem{Pinciu2015}
V. Pinciu.
\newblock Guarding polyominoes, polycubes and polyhypercubes.
\newblock \emph{Electronic Notes in Discrete Mathematics}, 49, 159--166, 2015.
\newblock \url{https://doi.org/10.1016/j.endm.2015.06.024}
\newblock Open FWCG version:
\url{https://cse.buffalo.edu/fwcg2015/assets/pdf/FWCG_2015_paper_14.pdf}

\bibitem{SchuchardtHecker1995}
D. Schuchardt and H.-D. Hecker.
\newblock Two NP-hard art-gallery problems for orthogonal polygons.
\newblock \emph{Mathematical Logic Quarterly}, 41(2), 261--267, 1995.
\newblock \url{https://doi.org/10.1002/malq.19950410212}

\bibitem{Shermer1992}
T. C. Shermer.
\newblock Recent results in art galleries.
\newblock \emph{Proceedings of the IEEE}, 80(9), 1384--1399, 1992.
\newblock \url{https://doi.org/10.1109/5.163407}

\bibitem{Urrutia2000}
J. Urrutia.
\newblock Art gallery and illumination problems.
\newblock In J.-R. Sack and J. Urrutia, editors, \emph{Handbook of Computational Geometry}, 973--1027. Elsevier, 2000.
\newblock \url{https://doi.org/10.1016/B978-044482537-7/50023-1}

\bibitem{UrrutiaOpenProblems}
J. Urrutia.
\newblock Illumination of polygons: open problems.
\newblock \url{https://www.matem.unam.mx/~urrutia/openprob/Polygons/}

\bibitem{VigliettaThesis}
G. Viglietta.
\newblock Guarding and searching polyhedra.
\newblock Ph.D. thesis; arXiv:1211.2483, 2012.
\newblock \url{https://arxiv.org/abs/1211.2483}
\newblock \url{https://giovanniviglietta.com/papers/thesis.pdf}

\bibitem{VigliettaFace2014}
G. Viglietta.
\newblock Face-guarding polyhedra.
\newblock \emph{Computational Geometry: Theory and Applications}, 47(8), 833--846, 2014.
\newblock \url{https://doi.org/10.1016/j.comgeo.2014.04.009}

\bibitem{Viglietta2020}
G. Viglietta.
\newblock Optimally guarding 2-reflex orthogonal polyhedra by reflex edge guards.
\newblock \emph{Computational Geometry: Theory and Applications}, 86, 101589, 2020.
\newblock \url{https://doi.org/10.1016/j.comgeo.2019.101589}
\newblock \url{https://arxiv.org/abs/1708.05469}

\bibitem{WormanKeil2007}
C. Worman and J. M. Keil.
\newblock Polygon decomposition and the orthogonal art gallery problem.
\newblock \emph{International Journal of Computational Geometry \& Applications}, 17(2), 105--138, 2007.
\newblock \url{https://doi.org/10.1142/S0218195907002264}

\bibitem{Zylinski2006}
P. Zylinski.
\newblock Orthogonal art galleries with holes: A coloring proof of Aggarwal's theorem.
\newblock \emph{The Electronic Journal of Combinatorics}, 13(1), R20, 2006.
\newblock \url{https://doi.org/10.37236/1046}
\newblock \url{https://www.combinatorics.org/ojs/index.php/eljc/article/view/v13i1r20}

\end{thebibliography}

\end{document}
